\begin{enumerate}
\def\labelenumi{\arabic{enumi}.}
\item
  \begin{enumerate}
  \def\labelenumii{\arabic{enumii}.}
  \item
    \begin{enumerate}
    \def\labelenumiii{\arabic{enumiii}.}
    \item
      Dualidad en el álgebra de Boole.

      \begin{enumerate}
      \def\labelenumiv{\arabic{enumiv}.}
      \item
        Cambiando \emph{\textbf{simultáneamente}} todos ``0,1,+,*'' por
        (respectivamente) ``1,0,*,+'' (más en general, cambiando cada
        valor de un \textbf{símbolo constante} de B por su
        complementario, dejando las variables inalteradas, aunque la
        prueba inicial la vamos a hacer permitiendo sólo los símbolos
        ``1,0'') obtenemos a partir de una ex\-presión otra con el mismo
        valor de verdad que la primera. A medida que vaya\-mos
        introduciendo nuevos operadores veremos como extender esta
        propiedad de dualidad. Se demuestra por inducción sobre
        cualquier expresión bien formada, teniendo en cuenta que todos
        los axiomas son simétricos, duales.

        \begin{enumerate}
        \def\labelenumv{\arabic{enumv}.}
        \tightlist
        \item
          Prueba:
        \end{enumerate}
      \end{enumerate}
    \end{enumerate}
  \end{enumerate}
\end{enumerate}

\begin{quote}
Primero hay que construir un lenguaje adecuado para las expresiones
booleanas (tal como está sirve para hacer un programa calculadora de
expresiones de Boo\-le, aunque no programable, ni tampoco se pueden
definir variables nuevas).
\end{quote}

\begin{enumerate}
\def\labelenumi{\arabic{enumi}.}
\setcounter{enumi}{1}
\item
  \begin{enumerate}
  \def\labelenumii{\arabic{enumii}.}
  \item
    \begin{enumerate}
    \def\labelenumiii{\arabic{enumiii}.}
    \item
      \begin{enumerate}
      \def\labelenumiv{\arabic{enumiv}.}
      \item
        \[\begin{matrix}
        {{\lbrack\mathbf{\mathit{Constante}}\rbrack}{\mathbf{C}{: =}\mathbf{0}o\mathbf{1}}} \\
        {\mathbf{C}{: =}\mathbf{0}o\mathbf{1}o\mathbf{\Lambda}o\mathbf{\Lambda_{\mathrm{n}}}} \\
        {\mathit{dónde}{\mathbf{n} \in \mathbb{N}}{({\mathit{ha}\mathit{de}\mathit{ser}\mathit{un}\mathit{número}\mathit{concreto}})}y} \\
        {\mathbf{\Lambda},{\mathbf{\Lambda_{\mathrm{n}}} \in {{\{{\alpha,\beta,\gamma,\delta,\epsilon,\theta,\eta,\%,\iota,\chi,\kappa,\lambda,\mu,\nu,\psi,ο,\pi,\rho}\}} \cup}}} \\
        {{\cup {\{{\sigma,\tau,\varepsilon,\varphi,\varsigma,\vartheta,\xi,\upsilon,\zeta,\omega,\varrho,\varpi}\}}}.} \\
        {\mathit{El}\mathit{conjunto}\mathit{de}\mathit{símbolos}\mathit{constantes}\mathit{ha}\mathit{de}\mathit{tener}} \\
        {\mathit{un}\mathit{cardinal}\mathit{menor}o\mathit{igual}\mathit{que}\mathbf{\# B}.} \\
        {\mathit{Estos}\mathit{símbolos}\mathit{tienen}\mathit{un}\mathit{valor}\mathit{concreto}\mathit{booleano}} \\
        {y\mathit{no}\mathit{toman}\mathit{otro}\mathit{valor.}\mathit{Puede}\mathit{haber}\mathit{dos}} \\
        {\mathit{signos}\mathit{constantes}\mathit{distintos}\mathit{con}\mathit{el}\mathit{mismo}} \\
        {\mathit{valor}\mathit{booleano},\mathit{en}\mathit{cuyo}\mathit{caso}\mathit{podemos}\mathit{escribir}\mathit{la}} \\
        {\mathit{igualdad}\mathit{entre}\mathit{los}\mathit{símbolos}\mathit{constantes.}}
        \end{matrix}\]

        \[\begin{matrix}
        {{\lbrack\mathbf{\mathit{Variable}}\rbrack}\mathbf{V}{: =}\mathbf{X}o\mathbf{X_{\mathrm{n}}}} \\
        \mathit{dónde} \\
        {\mathbf{X},{\mathbf{X_{\mathrm{n}}} \in {{\{{a,b,c,d,e,f,g,h,i,j,k,l,m,n,o,p,q,r}\}} \cup}}} \\
        {{\cup {\{{s,t,u,v,w,x,y,z}\}}}\phantom{\{\}}} \\
        {y\mathit{dónde}\mathbf{n}{\in}\mathbf{\mathbb{N}}.}
        \end{matrix}\]

        \[\begin{matrix}
        {\mathit{Lo}\mathit{importante}\mathit{de}\mathit{la}\mathit{variable}\mathbf{X}\mathit{no}\mathit{es}\mathit{que}\mathit{varíe}\mathit{en}\mathit{todo}\mathbf{B},} \\
        {\mathit{sino}\mathit{que}\mathit{varíe}\mathit{en}o\mathit{recorra}{\mathbf{Q_{X}} \subseteq \mathbf{B}},\mathit{tal}\mathit{que}{\mathbf{Q_{X}} \neq \varnothing}.} \\
        {\mathit{Cuando}\mathit{decimos}\mathit{que}\mathit{el}\mathit{dual}\mathit{de}\mathit{una}\mathit{variable}\mathit{es}\mathit{la}\mathit{misma}\mathit{variable},} \\
        {\mathit{no}\mathit{quiere}\mathit{decir}\mathit{que}\mathit{esta}\mathit{variable}\mathit{dual}\mathit{recorra}\mathit{el}\mathit{mismo}\mathit{conjunto}} \\
        {\mathit{de}\mathit{booleanos}\mathit{que}\mathit{la}\mathit{inicial},\mathit{sino}\mathit{que}\mathit{si}\mathit{la}\mathit{inicial}\mathit{tenía}\mathit{elementos},} \\
        {\mathit{la}\mathit{dual}\mathit{también}\mathit{los}\mathit{tiene.}}
        \end{matrix}\]

        \[\begin{matrix}
        {{\lbrack\mathbf{\mathit{Literal}}\rbrack}\mathbf{\mathit{L.p}}{: =}\mathbf{V}o\mathbf{C}} \\
        {\mathbf{L}{: =}\mathbf{\mathit{L.p}}o\overline{\mathbf{\mathit{L.p}}}}
        \end{matrix}\]

        \[\begin{matrix}
        {{\lbrack\mathbf{\mathit{Término}\mathit{Simple}}\rbrack}\mathbf{\mathit{TS.p}}{: =}\mathbf{L}o\left( {{{\cdot \mathbf{L}} \cdot} + {{\cdot \mathbf{L}} \cdot}} \right)o\left( {{{{\cdot \mathbf{L}} \cdot \mathbf{\mathrm{\cdot}}} \cdot \mathbf{L}} \cdot} \right)o\overline{\mathbf{L}}} \\
        {\mathbf{\mathit{TS}}{: =}\mathbf{\mathit{TS.p}}o\mathbf{\overline{\mathit{TS.p}}}} \\
        {\mathit{dónde}\mathit{los}\mathit{signos} \cdot \mathit{débiles}\mathit{significan}\mathit{seguido}ó\mathit{concatenado.}}
        \end{matrix}\]

        \[\begin{matrix}
        {{\lbrack\mathbf{\mathit{Término}}\rbrack}\mathbf{\mathit{T.p}}{: =}\mathbf{\mathit{TS}}o\left( {{{{\cdot \mathbf{\mathit{T.p}}} \cdot \mathbf{+}} \cdot \mathbf{\mathit{TS}}} \cdot} \right)o\left( {{{{\cdot \mathbf{\mathit{T.p}}} \cdot \mathbf{\mathrm{\cdot}}} \cdot \mathbf{\mathit{TS}}} \cdot} \right)} \\
        {{\lbrack\mathbf{\mathit{Término}}\rbrack}\mathbf{T}{: =}\mathbf{\mathit{T.p}}o\overline{\mathbf{\mathit{T.p}}}} \\
        {\mathit{dónde}\mathit{los}\mathit{signos} \cdot \mathit{débiles}\mathit{significan}\mathit{seguido}ó\mathit{concatenado.}}
        \end{matrix}\]

        \[{\lbrack\mathbf{\mathit{Expresión}}\rbrack}\mathbf{E}{: =}\mathbf{T{=}T}o\mathbf{T{\neq}T}\]

        Esquema de demostración:

        \begin{enumerate}
        \def\labelenumv{\arabic{enumv}.}
        \item
          Hay que demostrar que para cualquier expresión \[E\]pode\-mos
          encontrar la dual \[\widetilde{E}\](hasta aquí es sólo
          formalizar las reglas de los cambios) y que el valor de verdad
          de \[E\] y de \[\widetilde{E}\] son idénticos. Para ello vamos
          a demos\-trarlo primero para igualdades entre literales y
          des\-pués por recursividad para una expresión cualquiera con
          un literal, y por último para igualdades entre términos
          cualquiera. El método de demostración será el de induc\-ción
          matemática sobre el número de términos simples en ambos
          términos de una expresión. Lo haremos indepen\-dientemente
          para el caso de la expresión afirmada y para el caso de la
          expresión negada. Primero (1) se de\-muestra cuando los
          términos son literales, después (2) se demuestra cuando los
          términos son términos simples. El siguiente paso (3) es
          asegurarlo cuando un término es un término simple y el otro un
          término operado con un término simple, y suponiendo que ya se
          da la duali\-dad sin el añadido (HI), demostrar que se da la
          duali\-dad con el añadidos. El siguiente paso (4) es suponer
          que si es demostrar cuando los dos son términos supo\-niendo
          dado el caso que un término igual o desigual a un término
          simple está dado. Por último (5) dado que (HI) implica (3) y
          (HI) implica (4), suponiendo (2) de\-mostramos el caso término
          operado término simple y tér\-mino operado término simple (son
          3 casos de operacio\-nes) dado (HI). Estamos aprovechando para
          la demostra\-ción la estructura de la gramática construida
          anterior\-mente. En definitiva tendremos los puntos (a.1.cc) ,
          (a.1.cv) , (a.1.vv) , (a.1.lit.neg) , (a.2.lit.lit) ,
          (a.2.cc+c) , (a.2.vc+c) , (a.2.vv+c) , (a.2.cc+v) , (a.2.vc+v)
          , (a.2.vv+v) , (a.2.cc*c) , (a.2.vc*c) , (a.2.vv*c) ,
          (a.2.cc*v) , (a.2.vc*v) , (a.2.vv*v) , (a.2.ts.neg) ,
          (a.3.HI.tts+ts) , (a.3.HI.tts*ts) , (a.3.HI.tts+ts) ,
          (a.4.HI.ttsts*) , (a.5.HI.tts+tts+) , (a.5.HI.tts*tts*) ,
          (a.5.HI.tts+tts*) , (a.5.HI.t.neg) , y a continuación los
          puntos dónde se prueba que el dual de una desigual\-dad es
          equivalente a la dual de la misma desigualdad. Habrá que ver
          previamente que ver cuál es el dual de un recorrido \textbf{Q}
          en \textbf{B}, que será el conjunto de los elementos de
          \textbf{B} que son inversos de los elementos de \textbf{Q}, ya
          que para cada constante o valor que tome una variable, el dual
          será el dual de esa constante o valor, esto es, el complemento
          de esa constante o valor. Para una variable \textbf{a}, el
          conjunto de valores que recorre es
          \textbf{Q}\textsubscript{\textbf{a}} . Si necesi\-tamos algún
          conjunto auxiliar lo llamaremos \textbf{A} y se de\-finirá
          antes de su uso con la idea que no sea demasiado engorrosa la
          escritura. Si un conjunto \textbf{C} es sin el 1 lo
          denotaremos por \textbf{C}\textsubscript{\textbf{*}} y si es
          sin el 0 lo denotaremos \textbf{C}\textsuperscript{\textbf{*}}
          y definimos
          \textbf{B}\textsubscript{\textbf{2}}\textbf{=\{0,1\}} y si un
          conjunto \textbf{C} es sin 0 y sin 1 lo denotamos por
          \textbf{C}\textsuperscript{\textbf{X}}. Otro elemento
          importante son los paréntesis, que permiten una no
          asociatividad generali\-zada entre las diversas operaciones,
          al no cambiar en absoluto en la expresión dual.
        \item
          Casos en que la expresión \[E\] es\[T_{1} = T_{2}\]afirmativa:

          \emph{\textbf{a.1.cc)}} \[\Lambda_{1} = \Lambda_{2}\]. La
          expresión dual es
          \[\overline{\Lambda_{1}} = \overline{\Lambda_{2}}\]. Existe
          equivalencia cla\-ra entre estas 2 expresiones duales.

          \emph{\textbf{a.1.cv)}} \[\Lambda = X\]. El conjunto
          \[Q_{X{(\Lambda)}} = \left\{ \Lambda \right\}\] , y así es no
          vacío. El conjunto dual será
          \[{\widetilde{Q}}_{X{(\Lambda)}} = \left\{ \overline{\Lambda} \right\}\],
          que es evidentemente no vacía. La expresión dual es
          \[\overline{\Lambda} = X\].

          \emph{\textbf{a.1.vv)}} \[X = Y\]. El conjunto \[Q_{X} = B\] ,
          y así es no vacío. El conjunto asociado a \[Y\] es dependiente
          de los valores de \[X\]:
          \[Q_{Y{(X)}} = {B \smallsetminus {\{{\mathit{value}{(X)}}\}}}\]dónde
          \[\#{\left( \left\{ {\mathit{value}{(X)}} \right\} \right) = 1}\].
          De aquí que el con\-junto de valores que recorre \[Y\]no es
          vacío, esto es, \[Q_{Y{(X)}} \neq \varnothing\]. La
          ex\-presión dual será \[\widetilde{X} = \widetilde{Y}\] que se
          sigue traduciendo como \[X = Y\], ya que los conjunto
          asociados duales que son
          \[{{{\widetilde{Q}}_{X} = \left\{ {{\overline{\alpha} \in B} \mid {\alpha \in B}} \right\}} = B} = Q_{X}\],
          y
          \[{\widetilde{Q}}_{Y{(X)}} = {B \smallsetminus \left\{ \overline{\mathit{value}{(X)}} \right\}}\]sigue
          siendo no vacía. La expresión dual es evi\-dentemente
          equivalente para cada valor de \[X\].

          \emph{\textbf{a.1.lit.neg)}} Cada literal negado
          \[\overline{L}\] es o una variable negada o una cons\-tante
          negada. Cada constante negada no es más que otra constante.
          Luego si el literal negado es una constante sigue valiendo
          idénticamente los puntos anteriores. Si el literal negado es
          una variable digamos \[\overline{X}\] convertimos
          \[Q_{X} \neq \varnothing\] en
          \[Q_{\overline{X}} = \left\{ {{\overline{x} \in B} \mid {x \in Q_{X}}} \right\}\]
          que es la misma expresión de \[{\widetilde{Q}}_{X}\], y por lo
          tanto
          \[\#{{(Q_{X})} = \#}{{(Q_{\overline{X}})} = \#}{({\widetilde{Q}}_{X})}\]y
          además
          \[\#{{({\widetilde{Q}}_{X})} = \#}{{({\widetilde{Q}}_{\overline{X}})} = \#}{{({\widetilde{\widetilde{Q}}}_{X})} = \#}{(Q_{X})}\].
          Luego podemos también ver que se trata un caso reducido
          completamente a los anteriores. Queda demostrado que
          \[{L_{1} = L_{2}}\Leftrightarrow{\widetilde{L_{1}} = \widetilde{L_{2}}}\]
          .

          \emph{\textbf{a.2.lit.lit)}} Si
          \[\mathit{TS}_{1} = \mathit{TS}_{2}\] es tal que
          \[\mathit{TS}_{1} = L_{1}\] y \[\mathit{TS}_{2} = L_{2}\] ya
          ha queda\-do demostrado en el punto (a.1).

          \emph{\textbf{a.2.cc+c)}} El caso
          \[{\Lambda_{1} + \Lambda_{2}} = \Lambda_{3}\] tiene como
          expresión dual
          \[{\overline{\Lambda_{1}} \cdot \overline{\Lambda_{2}}} = \overline{\Lambda_{3}}\]que
          viene justificado por las leyes de De Morgan: ambas
          expresiones son equivalentes.

          \emph{\textbf{a.2.vc+c)}} El caso \[{X + \Lambda} = \Gamma\]
          tiene como expresión dual
          \[{X \cdot \overline{\Lambda}} = \overline{\Gamma}\] que viene
          justificado por las leyes de De Morgan, con sólo asegurar que
          en am\-bas expresiones \[X\]recorre un conjunto no vacío.
          Existen algunas posibili\-dades que de inicio sea vacío el
          recorrido de \[X\], por ejemplo si \[\Gamma = 0\] y
          \[\Lambda = 1\]. En estos casos comprobaremos que
          \[\widetilde{X}\] recorre un conjunto \[\varnothing\].
          Supongamos que
          \[\forall{x \in B}{{x + \Lambda} \neq \Gamma}\]y supongamos
          que
          \[\exists{y \in B}{{y \cdot \overline{\Lambda}} = \overline{\Gamma}}\].
          Y a la inversa: si existe un valor de \[X\] que haga la
          ex\-presión verdadera entonces existe al menos otro valor de
          \[\widetilde{X}\] que hará a la expresión dual igualmente
          verdadera.

          Prueba cuando \[X\] no tiene solución en la expresión:

          \begin{enumerate}
          \def\labelenumvi{\arabic{enumvi}.}
          \tightlist
          \item
            \begin{enumerate}
            \def\labelenumvii{\arabic{enumvii}.}
            \tightlist
            \item
              \[\forall{x \in B}{{x + \Lambda} \neq \Gamma}\]\[\land\]\[\exists{y \in B}{{y \cdot \overline{\Lambda}} = \overline{\Gamma}}\].
            \item
              \[{y \cdot \overline{\Lambda}} = \overline{\Gamma}\]
            \item
              \[{\left( {y \cdot \overline{\Lambda}} \right) + \left( {x + \Lambda} \right)} = {\overline{\Gamma} + \left( {x + \Lambda} \right)}\]
            \item
              \[{\left( {y \cdot \Lambda} \right) + x} = {\overline{\Gamma} + \left( {x + \Lambda} \right)}\]
            \item
              \[{{\left( {y + \Lambda} \right) + \overline{\Gamma}} + x} = {\overline{\Gamma} + \left( {x + \Lambda} \right)}\]
            \item
              \[{y = 0} = \overline{\Gamma}\]
            \item
              \[\Gamma = 1\]
            \item
              Haciendo \[x = \overline{\Lambda}\] tenemos que 1. es
              falso.
            \item
              \[\forall{x \in B}{{x + \Lambda} \neq \Gamma}\]\[\Rightarrow\]\[\forall{x \in B}{{x \cdot \overline{\Lambda}} \neq \overline{\Gamma}}\].
            \end{enumerate}
          \end{enumerate}

          Prueba cuando \[X\] no tiene solución en la expresión dual:

          \begin{enumerate}
          \def\labelenumvi{\arabic{enumvi}.}
          \tightlist
          \item
            \begin{enumerate}
            \def\labelenumvii{\arabic{enumvii}.}
            \tightlist
            \item
              \[\exists{x \in B}{{x + \Lambda} = \Gamma}\]\[\land\]\[\forall{y \in B}{{y \cdot \overline{\Lambda}} \neq \overline{\Gamma}}\].
            \item
              \[{x + \Lambda} = \Gamma\]
            \item
              \[{\left( {y \cdot \overline{\Lambda}} \right) + \left( {x + \Lambda} \right)} = {\overline{\Gamma} + \left( {x + \Lambda} \right)}\]
            \item
              \[{\left( {x + \overline{\Lambda}} \right) \cdot y} = {\Gamma \cdot \left( {y \cdot \overline{\Lambda}} \right)}\]
            \item
              \[{\left( {x + \overline{\Lambda}} \right) \cdot {({y \cdot \Gamma})}} = {\Gamma \cdot \left( {y \cdot \overline{\Lambda}} \right)}\]
            \item
              \[1 = {x \cdot \overline{\Gamma}}\]
            \item
              \[\Gamma = 0\]
            \item
              Haciendo \[y = \Lambda\] tenemos que 1. es falso.
            \item
              \[\forall{x \in B}{{x \cdot \overline{\Lambda}} \neq \overline{\Gamma}}\]\[\Rightarrow\]\[\forall{x \in B}{{x + \Lambda} \neq \Gamma}\].
            \end{enumerate}
          \end{enumerate}

          Prueba cuando \[X\] tiene solución en la expresión:

          \begin{enumerate}
          \def\labelenumvi{\arabic{enumvi}.}
          \tightlist
          \item
            \begin{enumerate}
            \def\labelenumvii{\arabic{enumvii}.}
            \tightlist
            \item
              \[\exists{Q_{X} \neq \varnothing}{Q_{X} \subseteq B}\forall{\xi \in Q_{X}}{{\xi + \Lambda} = \Gamma}\]
            \item
              \[\exists{Q_{X} \neq \varnothing}{Q_{X} \subseteq B}\forall{\xi \in Q_{X}}{{\overline{\xi} \cdot \overline{\Lambda}} = \overline{\Gamma}}\]
            \item
              \[\exists{\widetilde{Q_{X}} \neq \varnothing}{\widetilde{Q_{X}} \subseteq B}\forall{\xi \in \widetilde{Q_{X}}}{{\xi \cdot \overline{\Lambda}} = \overline{\Gamma}}\]
            \item
              \[{X \cdot \overline{\Lambda}} = \overline{\Gamma}\]
            \item
              \[{X + \Lambda} = \Gamma\] \[\Rightarrow\]
              \[{X \cdot \overline{\Lambda}} = \overline{\Gamma}\]
            \end{enumerate}
          \end{enumerate}

          Prueba cuando \[X\] tiene solución en la expresión dual:

          \begin{enumerate}
          \def\labelenumvi{\arabic{enumvi}.}
          \tightlist
          \item
            \begin{enumerate}
            \def\labelenumvii{\arabic{enumvii}.}
            \tightlist
            \item
              \[\exists{Q_{X} \neq \varnothing}{Q_{X} \subseteq B}\forall{\xi \in Q_{X}}{{\xi \cdot \overline{\Lambda}} = \overline{\Gamma}}\]
            \item
              \[\exists{Q_{X} \neq \varnothing}{Q_{X} \subseteq B}\forall{\xi \in Q_{X}}{{\overline{\xi} + \Lambda} = \Gamma}\]
            \item
              \[\exists{\widetilde{Q_{X}} \neq \varnothing}{\widetilde{Q_{X}} \subseteq B}\forall{\xi \in \widetilde{Q_{X}}}{{\xi + \Lambda} = \Gamma}\]
            \item
              \[{X + \Lambda} = \Gamma\]
            \item
              \[{X \cdot \overline{\Lambda}} = \overline{\Gamma}\]
              \[\Rightarrow\] \[{X + \Lambda} = \Gamma\]
            \end{enumerate}
          \end{enumerate}

          \emph{\textbf{a.2.vv+c)}} \[{X + Y} = \Gamma\]y esta expresión
          va a ser verdadera. Tenemos que ver su dual
          \[{X + Y} = \overline{\Gamma}\]es también verdadero. Para esto
          vemos la forma que to\-man los conjuntos asociados a las
          variables \[X\text{e}Y\]. Sea
          \[Q_{X} = {\{{{\alpha \cdot \Gamma} \mid {\alpha \in B}}\}}\]
          y
          \[Q_{Y{(X)}} = {\{{{\overline{\mathit{value}{(X)}} \cdot \Gamma} \mid \mathit{value}{{(X)} \in Q_{X}}}\}}\].
          Así en la expresión dual
          \[\widetilde{Q_{X}} = {\{{{\overline{\alpha} + \overline{\Gamma}} \mid {\alpha \in B}}\}}\]
          y en
          \[Q_{Y{(X)}} = {\{{{{\mathit{value}{(X)}} + \overline{\Gamma}} \mid \mathit{value}{{(X)} \in \widetilde{Q_{X}}}}\}}\],
          no son vacíos ninguno de los con\-juntos asociados, como en la
          expresión no dual. De hecho a pares tienen el mismo cardinal.
          Por lo tanto si uno de ellos fuera vacío para una variable
          también lo sería la en expresión dual, y viceversa.

          \emph{\textbf{a.2.cc+v)}} Veamos el caso
          \[{\Gamma + \Lambda} = X\]. La variable solo puede tomar un
          va\-lor, pero este existe
          \[Q_{X} = \left\{ {\Gamma + \Lambda} \right\}\]. Luego
          \[Q_{X}\]no es vacío. La expresión dual
          \[{\overline{\Gamma} \cdot \overline{\Lambda}} = X\]es tal que
          \[X\]recorre
          \[{{\widetilde{Q}}_{X} = \left\{ {\overline{\alpha} \mid {\alpha \in Q_{X}}} \right\}} = \left\{ {\overline{\Gamma} \cdot \overline{\Lambda}} \right\}\],
          que es no vacío, y con el mismo cardinal que \[Q_{X}\]. El
          camino inverso es similar, probado ya que los cardinales son
          los mismos.

          \emph{\textbf{a.2.vc+v)}} Este caso corresponde a
          \[{X + \Gamma} = Y\]. Su expresión dual sería
          \[{X \cdot \overline{\Gamma}} = Y\]. Si la expresión es
          verdadera \[{Q_{X} = B} \neq \varnothing\]y
          \[{Q_{Y{(X)}} = \left\{ {\mathit{value}{{(X)} + \Gamma}} \right\}} \neq \varnothing\].
          La expresión de los conjuntos asociados a las variables \[X\]
          e \[Y\] en la expresión dual sería
          \[{{\widetilde{Q}}_{X} = B} \neq \varnothing\], y
          \[{{\widetilde{Q}}_{Y{(X)}} = \left\{ {{\mathit{value}{(X)}} \cdot \overline{\Gamma}} \right\}} \neq \varnothing\].
          Al tener los mismos cardinales queda demos\-trado lo que
          queríamos demostrar para el caso afirmativo. Esta expresión y
          su dual son siempre verdaderas.

          \emph{\textbf{a.2.vv+v)}} El caso \[{X + Y} = Z\] tiene como
          dual \[{X \cdot Y} = Z\]. Solo tenemos que ver los conjuntos
          asociados \[{Q_{X} = B} \neq \varnothing\] ,
          \[{Q_{Y} = B} \neq \varnothing\] ,
          \[{Q_{Z{({X,Y})}} = \left\{ {{\mathit{value}{(X)}} + {\mathit{value}{(Y)}}} \right\}} \neq \varnothing\]
          , \[{{\widetilde{Q}}_{X} = B} \neq \varnothing\] ,
          \[{{\widetilde{Q}}_{Y} = B} \neq \varnothing\] ,
          \[{{{\widetilde{Q}}_{Z{({X,Y})}} = {\overline{\mathit{value}{(X)}} \cdot \overline{\mathit{value}{(Y)}}}} = \overline{{\mathit{value}{(X)}}+\overline{\mathit{value}{(Y)}}}} \neq \varnothing\]
          . Los cardinales de ambos jue\-gos de conjuntos asociados, los
          del original y los del dual son iguales. Esta expresión y su
          dual son siempre verdaderas.

          \emph{\textbf{a.2.cc*c)}} El caso
          \[{\Lambda_{1} \cdot \Lambda_{2}} = \Lambda_{3}\] tiene como
          expresión dual
          \[{\overline{\Lambda_{1}} + \overline{\Lambda_{2}}} = \overline{\Lambda_{3}}\]que
          viene justificado por las leyes de De Morgan: ambas
          expresiones son equivalentes.

          \emph{\textbf{a.2.vc*c)}} El caso
          \[{X \cdot \Lambda} = \Gamma\] tiene como expresión dual
          \[{X + \overline{\Lambda}} = \overline{\Gamma}\] que viene
          justificado por las leyes de De Morgan, con sólo asegurar que
          en am\-bas expresiones \[X\]recorre un conjunto no vacío.
          Existen algunas posibili\-dades que de inicio sea vacío el
          recorrido de \[X\], por ejemplo si \[\Gamma = 1\] y
          \[\Lambda = 0\]. En estos casos comprobaremos que
          \[\widetilde{X}\] recorre un conjunto \[\varnothing\].
          Supongamos que
          \[\forall{x \in B}{{x \cdot \Lambda} \neq \Gamma}\]y
          supongamos que
          \[\exists{y \in B}{{y + \overline{\Lambda}} = \overline{\Gamma}}\].
          Y a la inversa: si existe un valor de \[X\] que haga la
          expresión verdadera entonces existe al menos otro valor de
          \[\widetilde{X}\] que hará a la expresión dual igualmente
          verdadera.

          Prueba cuando \[X\] no tiene solución en la expresión:

          \begin{enumerate}
          \def\labelenumvi{\arabic{enumvi}.}
          \tightlist
          \item
            \begin{enumerate}
            \def\labelenumvii{\arabic{enumvii}.}
            \tightlist
            \item
              \[\forall{x \in B}{{x \cdot \Lambda} \neq \Gamma}\]\[\land\]\[\exists{y \in B}{{y + \overline{\Lambda}} = \overline{\Gamma}}\]
            \item
              \[{y + \overline{\Lambda}} = \overline{\Gamma}\]
            \item
              \[{\left( {y + \overline{\Lambda}} \right) + \left( {x \cdot \Lambda} \right)} = {\overline{\Gamma} + \left( {x \cdot \Lambda} \right)}\]
            \item
              \[{\left( {y + \overline{\Lambda}} \right) + x} = {\overline{\Gamma} + \left( {x \cdot \Lambda} \right)}\]
            \item
              \[{{\left( {y + \overline{\Lambda}} \right) + \overline{\Gamma}} + x} = {\overline{\Gamma} + \left( {x \cdot \Lambda} \right)}\]
            \item
              \[{{\left( {y \cdot \Lambda} \right) + {\overline{\Gamma} \cdot \Lambda}} + {x \cdot \Lambda}} = {{\overline{\Gamma} \cdot \Lambda} + \left( {x \cdot \Lambda} \right)}\]
            \item
              \[{y \cdot \Lambda} = 0\]
            \item
              \[y = \overline{\Lambda}\] y de ahí \[\Gamma = \Lambda\]
            \item
              Haciendo \[x = 1\] tenemos que 1. es falso.
            \item
              \[\forall{x \in B}{{x \cdot \Lambda} \neq \Gamma}\]\[\Rightarrow\]\[\forall{y \in B}{{y + \overline{\Lambda}} \neq \overline{\Gamma}}\]
            \end{enumerate}
          \end{enumerate}

          Prueba cuando \[X\] no tiene solución en la expresión dual:

          \begin{enumerate}
          \def\labelenumvi{\arabic{enumvi}.}
          \tightlist
          \item
            \begin{enumerate}
            \def\labelenumvii{\arabic{enumvii}.}
            \tightlist
            \item
              \[\forall{x \in B}{{x + \overline{\Lambda}} \neq \overline{\Gamma}}\]\[\land\]\[\exists{y \in B}{{y \cdot \Lambda} = \Gamma}\]
            \item
              \[{y \cdot \Lambda} = \Gamma\]
            \item
              \[{{{({y \cdot \Lambda})} + x} + \overline{\Lambda}} = {{\Gamma + x} + \overline{\Lambda}}\]
            \item
              \[{{y + \Lambda} + x} = {{\Gamma + x} + \overline{\Lambda}}\]
            \item
              \[{{{\Gamma + y} + \Lambda} + x} = {{\Gamma + x} + \overline{\Lambda}}\]
            \item
              \[y = 0\]
            \item
              \[\Gamma = 0\]
            \item
              \[\forall{x \in B}{{{x + \overline{\Lambda}} \neq \overline{\Gamma}} = 1}\]
            \item
              Haciendo \[x = \Gamma\] tenemos que 8. es falso y 1. es
              falso.
            \item
              \[\forall{x \in B}{{x + \overline{\Lambda}} \neq \overline{\Gamma}}\]\[\Rightarrow\]\[\forall{y \in B}{{y \cdot \Lambda} \neq \Gamma}\].
            \end{enumerate}
          \end{enumerate}

          Prueba cuando \[X\] tiene solución en la expresión:

          \begin{enumerate}
          \def\labelenumvi{\arabic{enumvi}.}
          \tightlist
          \item
            \begin{enumerate}
            \def\labelenumvii{\arabic{enumvii}.}
            \tightlist
            \item
              \[\exists{Q_{X} \neq \varnothing}{Q_{X} \subseteq B}\forall{\xi \in Q_{X}}{{\xi \cdot \Lambda} = \Gamma}\].
            \item
              \[\exists{Q_{X} \neq \varnothing}{Q_{X} \subseteq B}\forall{\xi \in Q_{X}}{{\overline{\xi} + \overline{\Lambda}} = \overline{\Gamma}}\]
            \item
              \[\exists{\widetilde{Q_{X}} \neq \varnothing}{\widetilde{Q_{X}} \subseteq B}\forall{\xi \in \widetilde{Q_{X}}}{{\xi + \overline{\Lambda}} = \overline{\Gamma}}\]
            \item
              \[{X \cdot \overline{\Lambda}} = \overline{\Gamma}\]
            \item
              \[{X + \Lambda} = \Gamma\] \[\Rightarrow\]
              \[{X \cdot \overline{\Lambda}} = \overline{\Gamma}\]
            \end{enumerate}
          \end{enumerate}

          Prueba cuando \[X\] tiene solución en la expresión dual:

          \begin{enumerate}
          \def\labelenumvi{\arabic{enumvi}.}
          \tightlist
          \item
            \begin{enumerate}
            \def\labelenumvii{\arabic{enumvii}.}
            \tightlist
            \item
              \[\exists{Q_{X} \neq \varnothing}{Q_{X} \subseteq B}\forall{\xi \in Q_{X}}{{\xi + \overline{\Lambda}} = \overline{\Gamma}}\]
            \item
              \[\exists{Q_{X} \neq \varnothing}{Q_{X} \subseteq B}\forall{\xi \in Q_{X}}{{\overline{\xi} \cdot \Lambda} = \Gamma}\]
            \item
              \[\exists{\widetilde{Q_{X}} \neq \varnothing}{\widetilde{Q_{X}} \subseteq B}\forall{\xi \in \widetilde{Q_{X}}}{{\xi \cdot \Lambda} = \Gamma}\]
            \item
              \[{X \cdot \Lambda} = \Gamma\]
            \item
              \[{X + \overline{\Lambda}} = \overline{\Gamma}\]
              \[\Rightarrow\] \[{X \cdot \Lambda} = \Gamma\]
            \end{enumerate}
          \end{enumerate}

          \emph{\textbf{a.2.cc*v)}} Veamos el caso
          \[{\Gamma + \Lambda} = X\]. La variable solo puede tomar un
          va\-lor, pero este existe
          \[Q_{X} = \left\{ {\Gamma + \Lambda} \right\}\]. Luego
          \[Q_{X}\]no es vacío. La expresión dual
          \[{\overline{\Gamma} \cdot \overline{\Lambda}} = X\]es tal que
          \[X\]recorre
          \[{{\widetilde{Q}}_{X} = \left\{ {\overline{\alpha} \mid {\alpha \in Q_{X}}} \right\}} = \left\{ {\overline{\Gamma} \cdot \overline{\Lambda}} \right\}\],
          que es no vacío, y con el mismo cardinal que \[Q_{X}\]. El
          camino inverso es similar, probado ya que los cardinales son
          los mismos.

          \emph{\textbf{a.2.vc*v)}} Este caso corresponde a
          \[{X + \Gamma} = Y\]. Su expresión dual sería
          \[{X \cdot \overline{\Gamma}} = Y\]. Si la expresión es
          verdadera \[{Q_{X} = B} \neq \varnothing\]y
          \[{Q_{Y{(X)}} = \left\{ {\mathit{value}{{(X)} + \Gamma}} \right\}} \neq \varnothing\].
          La expresión de los conjuntos asociados a las variables \[X\]
          e \[Y\] en la expresión dual sería
          \[{{\widetilde{Q}}_{X} = B} \neq \varnothing\], y
          \[{{\widetilde{Q}}_{Y{(X)}} = \left\{ {{\mathit{value}{(X)}} \cdot \overline{\Gamma}} \right\}} \neq \varnothing\].
          Al tener los mismos cardinales queda demos\-trado lo que
          queríamos demostrar para el caso afirmativo. Esta expresión y
          su dual son siempre verdaderas.

          a.2.vv*c)\[{X \cdot Y} = \Gamma\]y esta expresión va a ser
          verdadera. Tenemos que ver su dual
          \[{X \cdot Y} = \overline{\Gamma}\]es también verdadero. Para
          esto vemos la forma que toman los conjuntos asociados a las
          variables \[X\text{e}Y\]. Sea
          \[Q_{X} = {\{{{\alpha + \Gamma} \mid {\alpha \in B}}\}}\] y
          \[Q_{Y{(X)}} = {\{{{\overline{\mathit{value}{(X)}} + \Gamma} \mid \mathit{value}{{(X)} \in Q_{X}}}\}}\].
          Así en la expresión dual
          \[\widetilde{Q_{X}} = {\{{{\overline{\alpha} \cdot \overline{\Gamma}} \mid {\alpha \in B}}\}}\]
          y en
          \[Q_{Y{(X)}} = {\{{{{\mathit{value}{(X)}} \cdot \overline{\Gamma}} \mid \mathit{value}{{(X)} \in \widetilde{Q_{X}}}}\}}\],
          no son vacíos ninguno de los conjuntos asociados, como en la
          expresión no dual. De hecho a pares tienen el mismo cardinal.
          Por lo tanto si uno de ellos fuera va\-cío para una variable
          también lo sería la en expresión dual, y viceversa.

          \emph{\textbf{a.2.vv*v)}} El caso \[{X \cdot Y} = Z\] tiene
          como dual \[{X + Y} = Z\]. Solo tenemos que ver los conjuntos
          asociados \[{Q_{X} = B} \neq \varnothing\] ,
          \[{Q_{Y} = B} \neq \varnothing\] ,
          \[{Q_{Z{({X,Y})}} = \left\{ {{\mathit{value}{(X)}} \cdot {\mathit{value}{(Y)}}} \right\}} \neq \varnothing\]
          , \[{{\widetilde{Q}}_{X} = B} \neq \varnothing\] ,
          \[{{\widetilde{Q}}_{Y} = B} \neq \varnothing\] ,
          \[{{{\widetilde{Q}}_{Z{({X,Y})}} = \left\{ {\overline{\mathit{value}{(X)}} + \overline{\mathit{value}{(Y)}}} \right\}} = \left\{ \overline{\left( {{\mathit{value}{(X)}}\cdot{\mathit{value}{(Y)}}} \right)} \right\}} \neq \varnothing\]
          . Los cardi\-nales de ambos juegos de conjuntos asociados, los
          del original y los del dual son iguales. Esta expresión y su
          dual son siempre verdaderas.

          \emph{\textbf{a.2.ts.neg)}} El caso
          \[\mathit{TS} = \overline{\mathit{TS.p}}\], queda reducido a
          los casos anteriores ya que \[\mathit{TS.p} = L\] ,
          \[\mathit{TS.p} = \left( {L + L} \right)\] ,
          \[\mathit{TS.p} = \left( {L \cdot L} \right)\], que tienen
          como térmi\-nos duales a
          \[{\mathit{TS.p} = \widetilde{L}} = L\] ,
          \[{\mathit{TS.p} = \left( {\widetilde{L} \cdot \widetilde{L}} \right)} = \left( {L \cdot L} \right)\]
          ,
          \[{\mathit{TS.p} = \left( {\widetilde{L} + \widetilde{L}} \right)} = \left( {L + L} \right)\]
          , que son casos todos ellos anteriores.

          \emph{\textbf{a.3.HI.tts+ts)}} La expresión sería una del tipo
          \[{T + \mathit{TS}_{1}} = \mathit{TS}_{2}\] ,
          \[T{{{' + \mathit{TS}_{0}} + \mathit{TS}_{1}} = \mathit{TS}_{2}}\],
          dónde \[T'\]suponemos (HI) que los valores de ver\-dad de
          \[T_{1}{' = T_{2}}'\] son idénticos a
          \[\widetilde{T_{1}}{' = \widetilde{T_{2}}}'\]. Sólo tenemos
          que de\-mostrar que
          \[T_{1}{{{' + \mathit{TS}_{0}} + \mathit{TS}_{1}} = T_{2}}'\].
          Ahora bien \[\mathit{TS}_{0} + \mathit{TS}_{1}\] podemos
          re\-ducirlo finalmente a un literal por evaluación,
          concretamente a una variable \[X\] o a una constante
          \[\Gamma\]. Así nos queda \[T_{1}{{' + L} = T_{2}}{' + 0}\]que
          entra dentro de la \[\mathit{HI}\]. Así queda demostrado este
          caso.

          \emph{\textbf{a.3.HI.tts*ts)}} La expresión sería una del tipo
          \[{T \cdot \mathit{TS}_{1}} = \mathit{TS}_{2}\] ,
          \[T{{{' \cdot \mathit{TS}_{0}} \cdot \mathit{TS}_{1}} = \mathit{TS}_{2}}\],
          dónde \[T'\]suponemos (HI) que los valores de ver\-dad de
          \[T_{1}{' = T_{2}}'\] son idénticos a
          \[\widetilde{T_{1}}{' = \widetilde{T_{2}}}'\]. Sólo tenemos
          que de\-mostrar que
          \[T_{1}{{{' \cdot \mathit{TS}_{0}} \cdot \mathit{TS}_{1}} = T_{2}}'\].
          Ahora bien \[\mathit{TS}_{0} \cdot \mathit{TS}_{1}\] podemos
          reducir\-lo finalmente a un literal por evaluación,
          concretamente a una variable \[X\] o a una constante
          \[\Gamma\]. Así nos queda
          \[T_{1}{{' \cdot L} = T_{2}}{' \cdot 1}\]que entra dentro de
          la \[\mathit{HI}\]. Así queda demostrado este caso.

          \emph{\textbf{a.5.HI.tts+tts+)}} La expresión sería una del
          tipo \[{T_{1} + \mathit{TS}_{2}} = {T_{2} + \mathit{TS}_{3}}\]
          ,
          \[T_{1}{{{' + \mathit{TS}_{0}} + \mathit{TS}_{2}} = T_{2}}{{' + \mathit{TS}_{1}} + \mathit{TS}_{3}}\],
          dónde para \[T'\] suponemos (HI), que los valores de verdad de
          \[T_{1}{{' + \mathit{TS}_{2}} = T_{2}}{' + \mathit{TS}_{3}}\]
          son idénticos a los de
          \[\widetilde{T_{1}}{{' \cdot \widetilde{\mathit{TS}_{2}}} = \widetilde{T_{2}}}{' \cdot \widetilde{\mathit{TS}_{3}}}\]
          y viceversa. Ahora bien \[\mathit{TS}_{0} + \mathit{TS}_{2}\]
          podemos redu\-cirlo finalmente a un literal por evaluación,
          concretamente a una variable \[X\] o a una constante
          \[\Gamma\], e idénticamente para
          \[\mathit{TS}_{1} + \mathit{TS}_{3}\]. Así nos queda
          \[T_{1}{{' + L_{1}} = T_{2}}{' + L_{2}}\]que entra dentro de
          la \[\mathit{HI}\]. Así queda demos\-trado este caso.

          \emph{\textbf{a.5.HI.tts*tts*)}} La expresión sería una del
          tipo
          \[{T_{1} \cdot \mathit{TS}_{2}} = {T_{2} \cdot \mathit{TS}_{3}}\]
          ,
          \[T_{1}{{{' \cdot \mathit{TS}_{0}} \cdot \mathit{TS}_{2}} = T_{2}}{{' \cdot \mathit{TS}_{1}} \cdot \mathit{TS}_{3}}\],
          dónde para \[T'\] suponemos (HI), que los valores de verdad de
          \[T_{1}{{' \cdot \mathit{TS}_{2}} = T_{2}}{' \cdot \mathit{TS}_{3}}\]
          son idénticos a los de
          \[\widetilde{T_{1}}{{' + \widetilde{\mathit{TS}_{2}}} = \widetilde{T_{2}}}{' + \widetilde{\mathit{TS}_{3}}}\]
          y viceversa. Ahora bien
          \[\mathit{TS}_{0} \cdot \mathit{TS}_{2}\] podemos redu\-cirlo
          finalmente a un literal por evaluación, concretamente a una
          variable \[X\] o a una constante \[\Gamma\], e idénticamente
          para \[\mathit{TS}_{1} \cdot \mathit{TS}_{3}\]. Así nos
          que\-da \[T_{1}{{' \cdot L_{1}} = T_{2}}{' \cdot L_{2}}\]que
          entra dentro de la \[\mathit{HI}\]. Así queda demostrado este
          caso.

          \emph{\textbf{a.5.HI.tts+tts*)}} La expresión sería una del
          tipo
          \[{T_{1} \cdot \mathit{TS}_{2}} = {T_{2} \cdot \mathit{TS}_{3}}\]
          ,
          \[{\left( {T_{1}{' + \mathit{TS}_{0}}} \right) + \mathit{TS}_{2}} = {\left( {T_{2}{' + \mathit{TS}_{1}}} \right) \cdot \mathit{TS}_{3}}\],
          dónde para \[T'\] suponemos (HI), que los valores de verdad de
          \[T_{1}{{' + \mathit{TS}_{2}} = T_{2}}{' + \mathit{TS}_{3}}\]
          son idénticos a los de
          \[\widetilde{T_{1}}{{' \cdot \widetilde{\mathit{TS}_{2}}} = \widetilde{T_{2}}}{' \cdot \widetilde{\mathit{TS}_{3}}}\]
          y viceversa. Ahora bien \[\mathit{TS}_{0} + \mathit{TS}_{2}\]
          podemos redu\-cirlo finalmente a un literal por evaluación,
          concretamente a una variable \[X\] o a una constante
          \[\Gamma\], e idénticamente para
          \[\mathit{TS}_{1} \cdot \mathit{TS}_{3}\].
          \[T_{2}{' \cdot \mathit{TS}_{3}}\]es lo mismo que una
          \[T_{3}\] cualquiera (de cualquier longitud). Así nos queda
          \[T_{1}{{' + L_{1}} = {T_{3} \cdot L_{2}}}\]que ya ha sido
          demostrado en \emph{a.5.tts+tts+)}. Así queda demostrado este
          caso.

          a.5.HI.t.neg) Cualquier T.p negado es susceptible de ser
          desarrollado (por las leyes de De Morgan) como un T.p y se
          aplica \emph{a.5.HI.*}.
        \item
          Casos en que la expresión \[E\]
          es\[T_{1} \neq T_{2}\]afirmativa:

          Estos casos se reducen a que alguna variable siempre es vacía
          tanto en la ex\-presión original como en la dual. Se reducen a
          los casos anteriores.
        \end{enumerate}
      \end{enumerate}
    \end{enumerate}
  \end{enumerate}
\end{enumerate}
